\section*{NeurIPS Paper Checklist}

\begin{enumerate}

\item {\bf Claims}
    \item[] Question: Do the main claims made in the abstract and introduction accurately reflect the paper's contributions and scope?
    \item[] Answer: \answerYes{}
    \item[] Justification: The abstract and Section~\ref{sec:introduction} state the main contributions: the Wiggle Framework decomposing judge robustness into Mechanical Consistency, Single-turn Conviction, and Multi-turn Persistence; empirical wiggle rates of 25--71\% under static pushback and 62--91\% under adaptive persuasion across 9 frontier judges and 14 judging tasks; the net-corrupting nature of pressure-induced flips; and the jury-majority-strength practical signal. Each claim is supported by experimental results in Sections~\ref{sec:results} and~\ref{sec:discussion}.

\item {\bf Limitations}
    \item[] Question: Does the paper discuss the limitations of the work performed by the authors?
    \item[] Answer: \answerYes{}
    \item[] Justification: Section~\ref{sec:limitations} explicitly enumerates the main limitations: borderline-item sampling (so wiggle rates may be inflated relative to a representative content distribution), absence of a comparable human-judge baseline under the same protocol, model-generated rather than human-authored arguments at L2--L6, and a fixed three-model L6 persuader pool. Appendix~\ref{sec:appendix-limitations} expands on these and adds caveats around domain coverage, sample sizes, and the scope of black-box analysis.

\item {\bf Theory assumptions and proofs}
    \item[] Question: For each theoretical result, does the paper provide the full set of assumptions and a complete (and correct) proof?
    \item[] Answer: \answerNA{}
    \item[] Justification: The paper is empirical and does not present theoretical results, theorems, or proofs.

    \item {\bf Experimental result reproducibility}
    \item[] Question: Does the paper fully disclose all the information needed to reproduce the main experimental results of the paper to the extent that it affects the main claims and/or conclusions of the paper (regardless of whether the code and data are provided or not)?
    \item[] Answer: \answerYes{}
    \item[] Justification: Section~\ref{sec:framework} defines the framework, the pressure ladder, and the wiggle and retention rate equations. Section~\ref{sec:dataset} specifies the 9 judge models, their reasoning configurations, the temperature setting (0 for all queries), and the 6 datasets with sample sizes. Appendix~\ref{sec:appendix-sampling} gives per-dataset sampling procedures and borderline filters. Appendix~\ref{sec:appendix-pressure-gen} documents the L1--L5 generation pipeline, the full L6 adaptive persuasion protocol (including the three-agent loop and observer reliability validation), and the jury baseline. Appendix~\ref{sec:appendix-prompts} reproduces every system prompt used (judge, challenge, observer, persuader). With API access to the same frontier models, the experiments are reproducible end-to-end.

\item {\bf Open access to data and code}
    \item[] Question: Does the paper provide open access to the data and code, with sufficient instructions to faithfully reproduce the main experimental results, as described in supplemental material?
    \item[] Answer: \answerNo{}
    \item[] Justification: The paper does not release code or generated data with the submission. However, all source datasets used (WildGuard, AEGIS, HH-RLHF, ToxiGen, MAGE, Paired Prompts) are publicly available and cited, and every system prompt, sampling procedure, and protocol step is documented verbatim in the appendices. Code release is planned for the camera-ready version pending internal review.

\item {\bf Experimental setting/details}
    \item[] Question: Does the paper specify all the training and test details (e.g., data splits, hyperparameters, how they were chosen, type of optimizer) necessary to understand the results?
    \item[] Answer: \answerYes{}
    \item[] Justification: No models are trained in this work; the experiments are entirely API-based inference against closed-source frontier judges. Section~\ref{sec:dataset} specifies model versions, reasoning configurations (default per-family: off for OpenAI GPT-5 family and Claude 4.6, on for Grok-4.1 Reasoning and both Gemini 3 models), and the temperature-0 setting used for all queries. Appendix~\ref{sec:appendix-sampling} specifies dataset splits, sample sizes (100 per safety/toxicity/AI-detection dataset, 50 prompt pairs per Paired Prompts rubric), and stratification.

\item {\bf Experiment statistical significance}
    \item[] Question: Does the paper report error bars suitably and correctly defined or other appropriate information about the statistical significance of the experiments?
    \item[] Answer: \answerYes{}
    \item[] Justification: Figure~\ref{fig:hero} reports 95\% bootstrap confidence intervals over per-example outcomes (1000 resamples). Figure~\ref{fig:datasets-and-survival} reports Binary and Likert results separately, with 95\% bootstrap CIs over per-model wiggle rates (left panels, capturing inter-model variability) and per-dataset retention rates (right panels, capturing cross-dataset variability), in both cases over 1000 resamples. The corrective-majority claim in Section~\ref{fo:structure} is backed by per-condition z-tests with $p$-values reported.

\item {\bf Experiments compute resources}
    \item[] Question: For each experiment, does the paper provide sufficient information on the computer resources (type of compute workers, memory, time of execution) needed to reproduce the experiments?
    \item[] Answer: \answerYes{}
    \item[] Justification: The experiments consist exclusively of API-based inference against closed-source frontier judge models (no model training, no GPU/CPU compute on our side). Total API call volume is bounded by approximately 9 models $\times$ 14 (dataset, rubric, scale) judging tasks $\times$ 6 pressure levels $\times$ $\sim$100 items $\times$ up to 10 turns for multi-turn levels, plus mechanical-consistency repeats (10 trials per condition $\times$ 3 conditions per (model, task)). End-to-end runtime is dominated by vendor API throughput limits.

\item {\bf Code of ethics}
    \item[] Question: Does the research conducted in the paper conform, in every respect, with the NeurIPS Code of Ethics \url{https://neurips.cc/public/EthicsGuidelines}?
    \item[] Answer: \answerYes{}
    \item[] Justification: The research uses publicly available datasets (with original publications cited) and queries publicly available LLM APIs. No new human-subject data collection was conducted. The datasets we judge against include safety, toxicity, and political-content categories; we discuss restrictive/permissive directional bias and the corrupting nature of pressure in Section~\ref{sec:results}, with no privacy or fairness violations introduced by our protocol.

\item {\bf Broader impacts}
    \item[] Question: Does the paper discuss both potential positive societal impacts and negative societal impacts of the work performed?
    \item[] Answer: \answerYes{}
    \item[] Justification: Section~\ref{sec:introduction} motivates the work in terms of LLM judges deployed in safety-critical and agentic settings. Section~\ref{sec:discussion-pressure} discusses negative implications --- the corrupting nature of pressure-induced flips means that multi-agent judging designs (debate, self-critique, ensemble disagreement) inherit a net-corrupting effect when they expose judges to L4--L6-style pressure. The conclusion frames the framework as a tool for the field to track verdict stability as judges expand into reward modeling and agentic evaluation.

\item {\bf Safeguards}
    \item[] Question: Does the paper describe safeguards that have been put in place for responsible release of data or models that have a high risk for misuse (e.g., pre-trained language models, image generators, or scraped datasets)?
    \item[] Answer: \answerNA{}
    \item[] Justification: The paper does not release pre-trained models, image generators, or scraped datasets. The L6 adaptive persuader uses prompt templates documented in Appendix~\ref{sec:appendix-prompts}; analogous adversarial prompting against LLM judges has been studied in prior work (cited in Section~\ref{sec:related_work}, ``LLM-on-LLM persuasion and debate''), and surfacing this attack surface is necessary for defenders to design judging pipelines that anticipate it.

\item {\bf Licenses for existing assets}
    \item[] Question: Are the creators or original owners of assets (e.g., code, data, models), used in the paper, properly credited and are the license and terms of use explicitly mentioned and properly respected?
    \item[] Answer: \answerYes{}
    \item[] Justification: All datasets used (WildGuard, AEGIS, HH-RLHF, ToxiGen, MAGE, Paired Prompts) are publicly available and their original publications are cited in Section~\ref{sec:dataset} and the bibliography. Each dataset is used in accordance with its original license and terms of use. The frontier judge models are accessed via their respective vendor APIs under standard usage terms.

\item {\bf New assets}
    \item[] Question: Are new assets introduced in the paper well documented and is the documentation provided alongside the assets?
    \item[] Answer: \answerNA{}
    \item[] Justification: The paper does not release new datasets, model weights, or code as part of this submission. The Wiggle Framework is described entirely in the paper and appendices, with all system prompts and protocols reproduced verbatim.

\item {\bf Crowdsourcing and research with human subjects}
    \item[] Question: For crowdsourcing experiments and research with human subjects, does the paper include the full text of instructions given to participants and screenshots, if applicable, as well as details about compensation (if any)?
    \item[] Answer: \answerNA{}
    \item[] Justification: The paper does not involve crowdsourcing or new research with human subjects. We use existing labeled datasets that were created by their original publishers with their own annotation procedures.

\item {\bf Institutional review board (IRB) approvals or equivalent for research with human subjects}
    \item[] Question: Does the paper describe potential risks incurred by study participants, whether such risks were disclosed to the subjects, and whether Institutional Review Board (IRB) approvals (or an equivalent approval/review based on the requirements of your country or institution) were obtained?
    \item[] Answer: \answerNA{}
    \item[] Justification: The paper does not involve human subjects research conducted by the authors; it analyzes verdicts produced by LLM judges on existing public datasets.

\item {\bf Declaration of LLM usage}
    \item[] Question: Does the paper describe the usage of LLMs if it is an important, original, or non-standard component of the core methods in this research? Note that if the LLM is used only for writing, editing, or formatting purposes and does \emph{not} impact the core methodology, scientific rigor, or originality of the research, declaration is not required.
    \item[] Answer: \answerYes{}
    \item[] Justification: LLMs are central to the core methodology in two roles beyond being the subjects of evaluation. First, three persuader models (GPT-5.4, Claude 4.6 Opus, Grok-4.1 Reasoning) generate the L2--L4 counterarguments and the L6 adaptive challenges; these are documented in Appendix~\ref{sec:appendix-pressure-gen}. Second, an observer model (GPT-5) extracts each judge's verdict from its free-form response in multi-turn settings; the observer's role and a manual reliability validation against 100 transcripts spanning all six datasets and three persuaders are documented in Appendix~\ref{sec:appendix-l6-protocol}. All system prompts for both roles are reproduced in Appendix~\ref{sec:appendix-prompts}.

\end{enumerate}
